\section{Conclusion}

We have presented BFS-HBM: a fully synthesizable, open-source SystemVerilog
microarchitecture for BFS traversal of CSR-encoded sparse graphs over
High-Bandwidth Memory. The design is motivated by the observation that graph
traversal achieves only $\eta \approx 40\%$ HBM bandwidth efficiency~\cite{fritz2026rmwwall},
not from insufficient memory bandwidth, but from row-activation scheduling
constraints intrinsic to DRAM timing. A purpose-built ASIC that co-designs the
memory access pattern, the AXI4 burst strategy, and the on-chip frontier management
can recover this efficiency and deliver deterministic sub-microsecond latency that
no general-purpose processor can match.

The key microarchitectural contributions are:
\begin{itemize}[leftmargin=*,itemsep=2pt]
\item A single-beat CSR row-pointer fetch that recovers both $\mathit{row\_ptr}[v]$
  and $\mathit{row\_ptr}[v+1]$ in one AXI4 transaction, eliminating a second HBM
  round trip per vertex.
\item An 8-neighbor beat buffer that amortizes HBM burst latency across 8 edge
  checks per AXI transaction.
\item A 2-cycle BRAM read-modify-write visited bitmap that uses 128~KB of on-chip
  SRAM to track visited state for $2^{20}$ vertices, incurring no external memory
  access for the bitmap.
\item A closed-form throughput model connecting the RTL microarchitecture to the
  JEDEC HBM4 timing efficiency result, showing convergence to $\eta \approx 40\%$
  at realistic graph degrees.
\end{itemize}

Cycle-accurate simulation at 250~MHz with a 20-cycle modeled HBM latency confirms
correct BFS ordering and level assignment for all vertices of the test graph in 514 cycles.

The immediate next step is FPGA validation on a Xilinx Alveo U280, which provides
HBM2e at the same AXI4 interface that BFS-HBM targets. Real HBM latency measurements
on power-law graphs will validate the throughput model and provide the timing data
needed to close the ASIC design budget.

\vspace{6pt}
\noindent\textbf{Open Source.}
All RTL, testbenches, synthesis scripts, and simulation infrastructure are available at
\url{https://github.com/quantumcelnav/bfs-hbm-engine} under the MIT license.
The companion platform repository (\textit{graph-silicon}) provides workload-specific
wrappers, an HBM arbiter for multi-engine configurations, and characterization
harnesses for HFT, graph database, and genomics verticals.

\vspace{6pt}
\noindent\textbf{Acknowledgments.} The author thanks the JEDEC JC-42 working group
for public access to the HBM2e and HBM4 specification documents.
